\chapter{Functions and Methods}\label{ch:functions}

Functions are Dylan's means of doing anything. This chapter describes the
two kinds of function, methods and generic functions; how a method's
parameter list says what arguments it takes; how a generic function
chooses which of its methods to run; and the operations that call
functions and examine them. The standard is \dylanbook{chapters 6
and 10}.

\section{Methods and generic functions}\label{sec:functions-overview}

Every function is an instance of the class
\code{<function>}\index{<function>@\code{<function>}}, and there are two
kinds.

A \term{method} (class \code{<method>}) is a piece of code with a parameter
list. The parameter list may give a \term{specializer} for each required
parameter, a class or a singleton that the corresponding argument must
match. A method can be called directly, like a procedure in Scheme or a
closure in Lisp; called with arguments that do not match its specializers,
it signals an error.

A \term{generic function} (class \code{<generic-function>}) holds a set of
methods. When it is called, it finds the methods whose specializers match
the arguments, the \term{applicable} methods, and runs the most specific of
them. This choice, \term{method dispatch}, is how Dylan does
object-oriented programming: behavior is attached to classes by defining
methods on generic functions, not by putting functions inside classes.

\begin{dylan}
(define-method double ((thing <number>))
  (+ thing thing))
(define-method double ((thing <sequence>))
  (concatenate thing thing))
(double 10)                   ; => 20
(double "the rain ")          ; => "the rain the rain "
(double '(a b c))             ; => (a b c a b c)
(double 'cup)                 ; => error: no applicable method
(instance? double <generic-function>)          ; => #t
(instance? (method (x) x) <method>)            ; => #t
(instance? double <function>)                  ; => #t
\end{dylan}

\noindent Functions are ordinary objects. They can be stored in variables
and data structures, passed as arguments and returned as results.

\section{Making methods: method}\label{sec:functions-method}

\index{method@\code{method}}
\begin{quote}
\code{(method \meta{parameter-list} \meta{form} \ldots)}
\end{quote}

\noindent The special form \code{method} makes and returns a method. When
the method is called, its parameters are bound to the arguments as lexical
variables and the forms, its \term{body}, are evaluated in order; the
method returns the values of the last form, or \code{\#f} if there are none.

\begin{dylan}
((method (a b) (+ a b)) 3 4)             ; => 7
(bind ((square (method (n) (* n n))))
  (square (square 3)))                   ; => 81
(sort '(3 1 2) test: (method (a b) (> a b)))   ; => (3 2 1)
\end{dylan}

\noindent A method made by \code{method} is an \term{anonymous method}. It
can be given to a function, as the test to \code{sort} above, or kept in a
variable. To give a group of local methods names that can call one another,
use \code{bind-methods} (Section~\ref{sec:control-bind-methods}).

\subsection{Closures}\label{sec:functions-closures}

A method can refer to the variables in scope where it was made, and it
keeps them for as long as it exists, even after the form that bound them
has returned. Such a method is a \term{closure}. The variables are shared,
not copied: assignments made through one closure are seen by every other
closure over the same variable.

\begin{dylan}
(define-method make-counter ()
  (bind ((count 0))
    (values (method () (set! count (+ count 1)))
            (method () count))))
(define-method demo ()
  (bind ((next current (make-counter)))
    (next)
    (next)
    (list (next) (current))))
(demo)                         ; => (3 3)
\end{dylan}

\noindent A method that returns a method is a common idiom. The method below
takes a function and returns one that calls it twice:

\begin{dylan}
(define-method twice ((f <function>))
  (method (#rest args)
    (apply f args)
    (apply f args)))
(define print-twice (twice (method (n) (print n))))
(print-twice 55)                 ; prints: 55
                               ; prints: 55
                               ; => 55
\end{dylan}

\section{Parameter lists}\label{sec:functions-parameters}

The complete syntax of a parameter list (\dylanbook{chapter 6}) is

\begin{quote}
\code{(\{\meta{required} | (\meta{required} \meta{specializer})\}*\\
\phantom{(}[\#next \meta{next}]\\
\phantom{(}[\#rest \meta{rest}]\\
\phantom{(}[\#key \{\meta{keyparam} | ([\meta{keyword}] \meta{keyparam} [\meta{default}])\}*])}
\end{quote}

\noindent The parts must appear in this order. \code{\#next}, \code{\#rest}
and \code{\#key} may also be written \code{!next}, \code{!rest} and
\code{!key} (Section~\ref{sec:syntax-parameter-tokens}).

\begin{center}
\begin{tabular}{@{}ll@{}}
\toprule
Parameter list & Accepts \\
\midrule
\code{(a b c)} & three required arguments \\
\code{((a <list>) b \#rest c)} & two required (the first a list), then any number \\
\code{(\#rest a)} & any number of arguments \\
\code{(a b \#key foo bar)} & two required, then keywords \code{foo:} and \code{bar:} \\
\code{(\#key bar baz bim)} & keywords only \\
\code{(\#rest all \#key fee fi)} & keyword arguments, all kept in \code{all} \\
\bottomrule
\end{tabular}
\end{center}

\subsection{Required parameters}

Required parameters come first, and a call must supply an argument for
each, in order. A required parameter is a name, or a list of a name and a
specializer: \code{(x <integer>)}. A parameter without a specializer
accepts any object, as if specialized on \code{<object>}. Specializers are
described in Section~\ref{sec:functions-singletons}.

Calling a method with too few or too many arguments, or with an argument
that does not match its specializer, is an error:

\begin{dylan}
((method ((s <string>)) s) "x")   ; => "x"
((method ((s <string>)) s) 3)     ; => error: 3 does not match <string>
((method (a #rest r) a))          ; => error: too few arguments
\end{dylan}

\subsection{Rest parameters}

After \code{\#rest} comes one name. When the method is called, the
arguments left after the required ones are put into a sequence, which
becomes the value of the rest parameter; it is empty if there are none.

\begin{dylan}
(define-method show-rest (a #rest b)
  (list a b))
(show-rest 10 20 30 40)        ; => (10 (20 30 40))
(show-rest 10)                 ; => (10 ())
\end{dylan}

\noindent In both implementations the sequence is a list, but a program
should rely only on its being a sequence.

\subsection{Keyword parameters}\label{sec:functions-keywords}

After \code{\#key} come the \term{keyword parameters}. Keyword arguments are
optional, are given after the required arguments, and may come in any
order; each is a keyword followed by a value. A keyword parameter is
written in one of four ways:

\begin{center}
\begin{tabular}{@{}ll@{}}
\toprule
Written & Meaning \\
\midrule
\code{\meta{name}} & keyword \code{\meta{name}:}, parameter \meta{name}, default \code{\#f} \\
\code{(\meta{name} \meta{default})} & keyword \code{\meta{name}:}, parameter \meta{name} \\
\code{(\meta{keyword} \meta{param})} & the keyword and the parameter named separately \\
\code{(\meta{keyword} \meta{param} \meta{default})} & the same, with a default \\
\bottomrule
\end{tabular}
\end{center}

\noindent The \meta{default} is an expression evaluated each time the method
is called without that keyword. It is evaluated in a scope that includes all
the parameters before it: the required parameters, the next-method and
rest parameters, and the earlier keyword parameters. Without a default, an
unsupplied keyword parameter is \code{\#f}.

\begin{dylan}
(define-method percolate (#key (brand 'maxwell-house)
                               (cups 4)
                               (strength (if (> cups 6) 'strong 'mild)))
  (list brand cups strength))
(percolate)                          ; => (maxwell-house 4 mild)
(percolate cups: 10 brand: 'folgers) ; => (folgers 10 strong)
\end{dylan}

\noindent Naming the keyword and the parameter separately lets a method take
a keyword such as \code{position:} without calling its parameter
\code{position}, which would shadow the module variable of that name, often
a generic function the method needs:

\begin{dylan}
(define-method place (widget #key (position: the-pos 0) (size: the-size))
  (list widget the-pos the-size))
(place 'w size: 3)                   ; => (w 0 3)
\end{dylan}

\noindent All the required arguments are taken before any keyword argument.
If a call supplies too few, keywords and values are taken as required
arguments:

\begin{dylan}
(define-method show-keys (req1 req2 #key foo)
  (list req1 req2 foo))
(show-keys 'one 'two foo: 'three)    ; => (one two three)
(show-keys foo: 'three)              ; => (foo: three #f)
\end{dylan}

\noindent The arguments after the required ones must be keyword/value pairs;
anything else is an error. If the same keyword is given more than once, the
leftmost occurrence is used.

\begin{dylan}
(define-method grow ((x <number>) #key (by 1)) (+ x by))
(grow 5 by: 2 by: 100)               ; => 7
(grow 5 'by 2)                       ; => error: not keyword/value pairs
\end{dylan}

\subsection{Combining \#rest and \#key}

When a parameter list has both, \code{\#rest} comes first. The rest
parameter then receives all the keyword arguments, keywords and values, as a
sequence, and the method accepts any keyword
(Section~\ref{sec:functions-keyword-checking}); the arguments must still be
keyword/value pairs.

\begin{dylan}
(define-method test (the-req #rest the-rest #key a b)
  (list the-req the-rest a b))
(test 1 a: 2 b: 3 c: 4)              ; => (1 (a: 2 b: 3 c: 4) 2 3)
\end{dylan}

\subsection{The next-method parameter}

After the required parameters may come \code{\#next} and a name. That
parameter receives the \term{next method}, used to call less specific
methods (Section~\ref{sec:functions-next-method}). \code{define-method} gives every
method a next-method parameter named \code{next-method} unless the parameter
list names one with \code{\#next}, so it is normally written only to choose
another name, or in a method made with \code{method} that will be added to a
generic function.

\subsection{Keyword argument checking}\label{sec:functions-keyword-checking}

A method accepts only the keywords its parameter list names after
\code{\#key}, unless it has a rest parameter, in which case it accepts any
keyword. A generic function accepts, for a given call, every keyword that
any of the methods applicable to that call accepts; so if any applicable
method has \code{\#rest}, any keyword is accepted. The check is made by the
generic function before it runs any method, and a method called through a
generic function or through \code{next-method} does not check again.
Calling a method or generic function with a keyword it does not accept is an
error.

\begin{dylan}
(define-method label ((x <object>) #key price)
  (list price x))
(define-method label ((x <sequence>) #key unit-price)
  (add x (* unit-price (size x))))
(define-method label ((x <list>) #rest info #key calories)
  (add x calories))
(label 'grape price: 189)                       ; => (189 grape)
(label 'grape price: 189 unit-price: 2)         ; => error
(label (vector 3 4 5) price: 189 unit-price: 2) ; => #(6 3 4 5)
(label (vector 3 4 5) calories: 9)              ; => error
(label (list 3 4 5) protein: 7 calories: 9)     ; => (9 3 4 5)
\end{dylan}

\noindent The first \code{label} method is the only one applicable to a
symbol, so only \code{price:} is accepted. A vector is a sequence, so both
\code{price:} and \code{unit-price:} are. For a list the third method
applies, and it has \code{\#rest}, so anything goes.

The two implementations differ on one point. A method whose parameter list
has \code{\#key} with no keywords after it accepts no keywords in the
compiler to C, as the rule above says; in the Common Lisp implementation it
accepts any.

\section{Defining methods: define-method}\label{sec:functions-define-method}

\index{define-method@\code{define-method}}
\begin{quote}
\code{(define-method \meta{name} \meta{parameter-list} \meta{form} \ldots)}
\end{quote}

\noindent \code{define-method} makes a method, as \code{method} would, and
adds it to the generic function in the module variable \meta{name}. If the
variable does not exist, it is created, holding a new generic function. If
the generic function already has a method with exactly the same
specializers, the new method replaces it. \meta{name} may be a setter name,
\code{(setter \meta{symbol})} (Section~\ref{sec:functions-setters}). At the
interactive prompt \code{define-method} returns the name.

%! lisp-only
\begin{dylan}
(define-method area ((r <integer>)) (* r r))
(area 3)                       ; => 9
(define-method area ((r <integer>)) (* 3 r r))
(area 3)                       ; => 27
\end{dylan}

\noindent That is the Common Lisp implementation, where each definition
takes effect when it is evaluated. In the compiler to C all of a program's
methods exist before it starts (Section~\ref{sec:syntax-when}), so the later
definition simply wins: both calls of \code{area} above return 27.

\noindent It is an error for \meta{name} to hold something other than a
generic function. The compiler to C also refuses to add methods to the
functions the manual calls methods rather than generic functions, such as
\code{id?}, \code{not}, \code{list} and \code{apply}, and to the
reflective functions it implements in C.

\subsection{Congruent parameter lists}

All the methods of a generic function, and its
\code{define-generic-function} if it has one, must have \term{congruent}
parameter lists: the same number of required parameters, and either all of
them have \code{\#rest} or \code{\#key} (or both), or none has either.
Adding a method that is not congruent is an error, reported when compiling
by the compiler to C.

\begin{center}
\begin{tabular}{@{}l@{}}
\toprule
Groups of congruent parameter lists \\
\midrule
\code{(a b)} \quad \code{(x y)} \\
\code{(a b \#key)} \quad \code{(x y \#key foo)} \quad
  \code{(a b \#rest l)} \quad \code{(a b \#rest l \#key quux)} \\
\code{(a b c)} \quad \code{(d e f)} \\
\bottomrule
\end{tabular}
\end{center}

\noindent When \code{define-method} creates the generic function, the
generic function takes the same number of required arguments as the method,
and, if the method has \code{\#rest} or \code{\#key}, a rest argument.
\code{function-arguments} (Section~\ref{sec:functions-reflection}) shows this:

\begin{dylan}
(define-method foo (a #key x) x)
(function-arguments foo)       ; => 1, #t, #f
\end{dylan}

\noindent A generic function with no required parameters can hold only one
method, since all its methods have the same (empty) specializers; defining
another replaces the first.

\section{Declaring generic functions: define-generic-function}\label{sec:functions-define-generic}

\index{define-generic-function@\code{define-generic-function}}
\begin{quote}
\code{(define-generic-function \meta{name} \meta{parameter-list} \meta{keyword} \meta{option} \ldots)}
\end{quote}

\noindent \code{define-generic-function} creates a generic function with no
methods and stores it in the module variable \meta{name}, which is made
read-only. It is never required, since \code{define-method} creates generic
functions as needed, but it documents the protocol in one place and fixes
the parameter list that every method must be congruent with.

The parameter list has the form

\begin{quote}
\code{(\meta{required} \ldots{} [\#rest \meta{rest}] [\#key \meta{keyword} \ldots])}
\end{quote}

\noindent Required parameters are names only, without specializers. The
keys may be written as symbols or keywords, and \code{\#key} may appear with
no keys at all. \code{\#key} means that the arguments after the required
ones must be keyword/value pairs; \code{\#rest} means any number of further
arguments. Either way, every method must have \code{\#rest} or \code{\#key}.
The meaning of giving both is unspecified.

\begin{dylan}
(define-generic-function resize (shape #key scale))
(define-method resize ((n <integer>) #key scale)
  (* n (or scale 1)))
(resize 4 scale: 3)            ; => 12
(function-arguments resize)    ; => 1, #f, (scale:)
\end{dylan}

\noindent When keys are listed, every method is meant to accept them, by
naming them after \code{\#key} or by having \code{\#rest}. The Common Lisp
implementation checks this when the method is added; the compiler to C does
not.

The options are keyword/value pairs reserved by the manual for future use.
Both implementations accept them and ignore them. In the compiler to C a
generic function may be declared only once.

\section{Specializers and singletons}\label{sec:functions-singletons}

A specializer is an expression, evaluated once when the method is made, not
on each call. It must evaluate to a class or to a singleton.

\begin{itemize}
\item A \term{class} specializer matches any \emph{general instance} of the
  class: a direct instance, or an instance of any subclass. The specializer
  \code{<number>} matches integers, ratios and floats.
\item A \term{singleton} specializer matches one object, compared with
  \code{id?}. \code{(singleton \meta{object})} makes one; so does
  \code{(make <singleton> object: \meta{object})}.
\end{itemize}

\begin{dylan}
(define-method factorial ((n <integer>))
  (* n (factorial (- n 1))))
(define-method factorial ((n (singleton 0)))
  1)
(factorial 5)                  ; => 120
(define-method measure ((x (singleton 'cup))) 'pint)
(measure 'cup)                 ; => pint
(instance? 3 (singleton 3))    ; => #t
\end{dylan}

\noindent \index{singleton@\code{singleton}}\code{singleton} is a function;
\code{<singleton>}\index{<singleton>@\code{<singleton>}} is the class of
singletons, and a singleton prints as \code{\{singleton 3\}}. Because
singletons compare their objects with \code{id?}, they suit symbols,
keywords, characters, small integers, classes and other objects whose
identity is what matters; a singleton of a string literal matches only that
very string object.

Singletons can also give an individual object slots of its own
(Section~\ref{sec:classes-define-slot}).

\begin{warnbox}
In the compiler to C every specializer of a method added to a generic
function must be known when compiling: a class or a variable whose value is
a class, or \code{(singleton \meta{x})} or \code{(make <singleton> object:
\meta{x})} where \meta{x} is a literal, a class, or a variable whose
\code{define} gives it a literal value. \code{(singleton (+ 1 2))}, and a
singleton of an object made at run time, are errors there, though the
Common Lisp implementation accepts them.
\end{warnbox}

\section{Method dispatch}\label{sec:functions-dispatch}

When a generic function is called, it
\begin{enumerate}
\item finds the \emph{applicable} methods: those whose every specializer
  matches the corresponding argument;
\item sorts them from most to least specific;
\item calls the most specific, passing it the next one as its next method.
\end{enumerate}
\noindent If no method is applicable, an error is signaled.

\subsection{Specificity}\label{sec:functions-specificity}

For one argument, a singleton specializer is more specific than any class,
and a class is more specific than its superclasses. So a method on
\code{<integer>} overrides one on \code{<number>} for integer arguments, and
a method on \code{(singleton 5)} overrides both for the argument 5.

\begin{dylan}
(define-method describe ((x <number>)) 'number)
(define-method describe ((x <integer>)) 'integer)
(define-method describe ((x (singleton 5))) 'five)
(describe 2.5)                 ; => number
(describe 2)                   ; => integer
(describe 5)                   ; => five
\end{dylan}

\noindent Methods that specialize more than one parameter are
\term{multimethods}, and dispatch considers all the specialized arguments.
Following \dylanbook{chapter 10}, both implementations sort the applicable
methods by comparing their specializer lists \emph{lexicographically}: the
first parameter decides, and a later parameter is looked at only when the
earlier specializers are the same.

\begin{dylan}
(define-method meet ((a <number>) (b <integer>)) 'number-integer)
(define-method meet ((a <integer>) (b <number>)) 'integer-number)
(meet 1 2)                     ; => integer-number
\end{dylan}

\noindent Both methods apply to \code{(meet 1 2)}. The first parameter's
specializers differ, and \code{<integer>} is more specific than
\code{<number>}, so the second method wins, though the first is more
specific in the second parameter.

\subsection{The class precedence list}\label{sec:functions-cpl}

When a class has several direct superclasses, two specializers can both be
superclasses of an argument's class without either being a subclass of the
other. They are then ordered by the argument's class's \term{class
precedence list}: the class, then all its superclasses, in an order that
keeps each class before its superclasses, keeps the direct superclasses in
the order the class's definition lists them, and keeps each simple chain of
superclasses together. \code{all-superclasses} returns this list
(Section~\ref{sec:classes-reflection}). The algorithm is the one CLOS uses, given
in design note 31.

\begin{dylan}
(define-class <life-form> (<object>))
(define-class <sentient> (<life-form>))
(define-class <bipedal> (<life-form>))
(define-class <intelligent> (<sentient>))
(define-class <humanoid> (<bipedal>))
(define-class <vulcan> (<intelligent> <humanoid>))
(define-class <human> (<humanoid> <intelligent>))
(all-superclasses <vulcan>)
  ; => ({the class <vulcan>} {the class <intelligent>}
  ;     {the class <sentient>} {the class <humanoid>}
  ;     {the class <bipedal>} {the class <life-form>}
  ;     {the class <object>})
(define-method psychoanalyze ((being <intelligent>)) 'intelligent)
(define-method psychoanalyze ((being <humanoid>)) 'humanoid)
(psychoanalyze (make <vulcan>)) ; => intelligent
(psychoanalyze (make <human>))  ; => humanoid
\end{dylan}

\noindent A \code{<vulcan>} lists \code{<intelligent>} first among its
superclasses, so the \code{<intelligent>} method is more specific for it; a
\code{<human>} lists \code{<humanoid>} first. With multimethods the same
rule is applied lexicographically:

%! cont
\begin{dylan}
(define-method superior ((a <intelligent>) (b <intelligent>)) 'smarter)
(define-method superior ((a <humanoid>) (b <humanoid>)) 'better-looking)
(superior (make <vulcan>) (make <human>))   ; => smarter
(superior (make <human>) (make <vulcan>))   ; => better-looking
\end{dylan}

\begin{notebox}
Design note 31 later replaced the lexicographic rule with one that compares
every argument position and calls the methods \emph{ambiguous}, an error,
when the positions disagree, as they do in the last two calls above. Neither
implementation follows that note: both use the 1992 rule, under which a call
whose applicable methods all have class or singleton specializers is never
ambiguous.
\end{notebox}

\subsection{Calling the next method}\label{sec:functions-next-method}

\index{next-method@\code{next-method}}
A method often wants to add to the behavior it overrides rather than
replace it. Its next-method parameter, normally named \code{next-method},
holds a function that calls the next most specific applicable method. Called
with no arguments, it passes along the arguments the current method
received; called with arguments, it passes those instead.

\begin{dylan}
(define-method greet ((x <object>)) (list 'object x))
(define-method greet ((x <integer>)) (cons 'integer (next-method)))
(define-method greet ((x (singleton 5))) (cons 'five (next-method 7)))
(greet "hi")                   ; => (object "hi")
(greet 3)                      ; => (integer object 3)
(greet 5)                      ; => (five integer object 7)
\end{dylan}

\noindent When new arguments are given, they must lead to the same
applicable methods, in the same order, as the original ones; otherwise the
result is undefined. (In \code{greet} above, 7 is still an integer and an
object.) Keyword arguments passed to \code{next-method} with changed values
are not checked.

When there is no next method, the parameter is \code{\#f}, so a method can
test it before calling it; calling \code{\#f} is an error.

\begin{dylan}
(define-method last-one ((x <object>)) next-method)
(last-one 1)                   ; => #f
(define-method chain ((x <number>) #next more) (list 'number more))
(define-method chain ((x <integer>) #next more) (list 'integer (more)))
(chain 1)                      ; => (integer (number #f))
\end{dylan}

\noindent When methods take different keyword arguments, giving each a rest
parameter lets a method pass all of them on, changed or not, with
\code{apply} (Section~\ref{sec:functions-apply}):

\begin{dylan}
(define-method grow ((x <number>) #rest options #key (by 1))
  (+ x by))
(define-method grow ((x <integer>) #rest options #key by)
  (apply next-method x by: 10 options))
(grow 5)                       ; => 15
(grow 5 by: 2)                 ; => 15
\end{dylan}

\noindent The second call shows the leftmost-keyword rule at work: the
\code{by: 10} the integer method puts first wins over the caller's
\code{by: 2}.

\section{Calling functions: apply}\label{sec:functions-apply}

\index{apply@\code{apply}}
\begin{quote}
\code{(apply \meta{function} \meta{argument} \ldots{} \meta{sequence})}
\end{quote}

\noindent \code{apply} calls \meta{function} with the given arguments
followed by the elements of \meta{sequence}, the last argument, and returns
its values.

\begin{dylan}
(apply + 1 '(2 3))             ; => 6
(apply max (vector 3 9 2))     ; => 9
(define math-functions (list + * -))
(apply (first math-functions) 1 2 '(3 4))   ; => 10
(apply list '())               ; => ()
\end{dylan}

\section{Setters as generic functions}\label{sec:functions-setters}

A setter is an ordinary function, normally a generic function, stored in a
setter variable \code{(setter \meta{name})}. Its arguments are those of the
getter followed by the new value. \code{define-class} defines setter methods
for slots, and a program can define its own with \code{define-method}, so
that \code{set!} works on a computed property as it does on a slot:

\begin{dylan}
(define-class <temperature> (<object>)
  (kelvin init-value: 0 init-keyword: kelvin:))
(define-method celsius ((t <temperature>))
  (- (kelvin t) 273))
(define-method (setter celsius) ((t <temperature>) c)
  (set! (kelvin t) (+ c 273)))
(define t1 (make <temperature>))
(set! (celsius t1) 25)         ; => 25
(kelvin t1)                    ; => 298
(celsius t1)                   ; => 25
\end{dylan}

\noindent Section~\ref{sec:control-extended-set} describes how \code{set!}
calls a setter. Later Dylan changed both the name, to \code{celsius-setter},
and the argument order, putting the new value first; neither change is part
of this language.

\section{Reflective operations}\label{sec:functions-reflection}

These functions examine functions and generic functions
(\dylanbook{chapter 10}). Ordinary programs seldom need them, but they are
useful in tools and for understanding dispatch.

\begin{description}
\item[\code{(function-arguments \meta{function})}]
  \index{function-arguments@\code{function-arguments}}
  returns three values: the number of required arguments, whether the
  function takes a rest argument, and a sequence of the keywords it
  accepts, or \code{\#f} if it has no \code{\#key}. The sequence is empty
  for \code{\#key} with no keywords.
\item[\code{(applicable-method? \meta{function} \meta{argument} \ldots)}]
  \index{applicable-method?@\code{applicable-method?}}
  returns true if \meta{function} is a method applicable to the arguments,
  or a generic function with such a method.
\item[\code{(sorted-applicable-methods \meta{generic-function} \meta{argument} \ldots)}]
  \index{sorted-applicable-methods@\code{sorted-applicable-methods}}
  returns a sequence of the methods applicable to the arguments, most
  specific first.
\item[\code{(find-method \meta{generic-function} \meta{specializers})}]
  \index{find-method@\code{find-method}}
  returns the method whose specializers are exactly the classes and
  singletons in the sequence \meta{specializers}, or \code{\#f}.
\item[\code{(generic-function-methods \meta{generic-function})}]
  returns a sequence of all its methods, in no particular order.
\item[\code{(method-specializers \meta{method})}]
  returns a sequence of a method's specializers, one for each required
  parameter.
\end{description}

\noindent None of the sequences these functions return should be modified.

\begin{dylan}
(function-arguments (method (a b #rest r) a))     ; => 2, #t, #f
(function-arguments (method (a #key x (y: z 3)) a)) ; => 1, #f, (x: y:)
(define-method kind ((x <number>)) 'number)
(define-method kind ((x <integer>)) 'integer)
(define-method kind ((x (singleton 5))) 'five)
(applicable-method? kind 1/2)                     ; => #t
(applicable-method? kind "x")                     ; => #f
(map (method (m) (m 5))
     (sorted-applicable-methods kind 5))          ; => (five integer number)
((find-method kind (list <integer>)) 7)           ; => integer
(find-method kind (list <string>))                ; => #f
(size (generic-function-methods kind))            ; => 3
(method-specializers (find-method kind (list <integer>)))
                                                  ; => ({the class <integer>})
\end{dylan}

\subsection{Changing generic functions}\label{sec:functions-changing}

The 1992 manual also has functions that change a generic function while the
program runs. \code{add-method}, \code{remove-method} and
\code{freeze-methods} exist only in the Common Lisp implementation; in the
compiler to C, whose generic functions are fixed when the program is
compiled, using them is an error when compiling. \code{make-read-only}
exists in both, but in the compiler to C it changes nothing, since no
generic function can change there anyway; it returns its argument.

\begin{description}
\item[\code{(add-method \meta{generic-function} \meta{method})}] adds the
  method, replacing any with the same specializers, and returns two values:
  the method, and the method it replaced or \code{\#f}.
\item[\code{(remove-method \meta{generic-function} \meta{method})}] removes
  the method and returns it.
\item[\code{(make-read-only \meta{generic-function})}] forbids adding or
  removing methods from then on.
\item[\code{(freeze-methods \meta{generic-function} \meta{specializer} \ldots)}]
  forbids adding, replacing or removing methods whose specializers are the
  given ones or more specific.
\end{description}

%! lisp-only
\begin{dylan}
(define-method double ((x <number>)) (+ x x))
(define m (method ((s <string>)) (concatenate s s)))
(add-method double m)          ; => {a method on double}, #f
(double "ab")                  ; => "abab"
(remove-method double m)       ; => {an anonymous method}
(double "ab")                  ; => error: no applicable method
\end{dylan}

\noindent A generic function can also be made with \code{make}, with the
init-keywords \code{required:} (the number of required arguments, which must
be given), \code{rest?:}, \code{key?:} and \code{debug-name:}. It starts with
no methods. In the compiler to C such a generic function can never have any,
since methods cannot be added at run time, so it is of little use there;
that implementation also ignores \code{rest?:} and \code{key?:}.

%! lisp-only
\begin{dylan}
(define expand
  (make <generic-function> required: 1 debug-name: 'expand))
(add-method expand (method ((n <integer>)) (* n 10)))
(expand 5)                     ; => 50
\end{dylan}

\section{Proper tail calls}\label{sec:functions-tail-calls}

A call is in \term{tail position} when its value is the value of the
enclosing method: the last form of a method's body, either branch of an
\code{if} in tail position, the last form of a \code{begin} or \code{bind}
in tail position, and so on. The compiler to C guarantees \term{proper tail
calls}: every call in tail position, to a known method, a generic function,
a closure or a function in a variable, runs in constant stack space. A
recursive method whose recursive calls are all tail calls is therefore a
loop, and two methods can call each other in tail position indefinitely.
The \code{for}, \code{while} and other loops of
Chapter~\ref{ch:control} are themselves compiled this way.

\begin{dylan}
(define-method my-even? ((n <integer>))
  (if (= n 0) #t (my-odd? (- n 1))))
(define-method my-odd? ((n <integer>))
  (if (= n 0) #f (my-even? (- n 1))))
(my-even? 10000000)            ; => #t
\end{dylan}

\noindent The Common Lisp implementation compiles Dylan to Common Lisp, and
SBCL eliminates tail calls in that code too, so the example above runs there
as well; but the Common Lisp standard does not require it, and the Common
Lisp implementation makes no guarantee. Chapter~\ref{ch:dylanc} describes how the
compiler to C implements tail calls, with the C compiler's \code{musttail}
or with trampolines.
